\documentclass[10pt,a4paper]{amsart}
\usepackage[margin=19mm]{geometry}
\usepackage{amsmath,amssymb,mathtools}
\usepackage[scr=rsfs,cal=euler]{mathalfa}
\usepackage{xfrac}
\usepackage[unicode,hidelinks,bookmarks=false]{hyperref}
\newcommand{\ie}{i.e.\ }
\newcommand{\cf}{cf.\ }
\newcommand{\resp}{respectively\ }
\newcommand{\Subsection}{\subsection}
\providecommand{\nobreakdash}{\nobreak\hbox{-}}
\setlength{\emergencystretch}{2em}
\multlinegap0pt
\usepackage{amssymb}
\newtheorem{proposition}{Proposition}
\title{Exact Hull Reformulation for Quadratically Constrained Generalized Disjunctive Programs}
\author{Sergey Gusev, David E. Bernal Neira}
\date{Source: arXiv:2508.16093v2 (CC BY 4.0)}
\begin{document}
\maketitle

\noindent\emph{Source and licence.} Exact Hull Reformulation for Quadratically Constrained Generalized Disjunctive Programs. Version arXiv:2508.16093v2 (CC BY 4.0), distributed by arXiv under the Creative Commons Attribution 4.0 International licence, http://creativecommons.org/licenses/by/4.0/, as recorded in the arXiv licence metadata for this exact version and shown on the abstract page https://arxiv.org/abs/2508.16093v2. Full licence notice: \texttt{licenses/arxiv-2508.16093-CC-BY-4.0.txt}.\par\smallskip

\noindent\textit{Editorial scope.} Counted unit: the article's proposition prop:conic-epi-equiv (source0.tex lines 677--706). The article's complete original proof of it is delivered with it (source0.tex lines 708--733, one \texttt{proof} environment of 26 lines). No mathematics is written, repaired or re-derived: the statement and the proof are the article's own lines, in the article's own notation, and no retained line is substituted or re-flowed.  Retained uncounted as the source-local context the proof needs: lines 617--628.  Nothing was omitted from any retained span, so the package carries no omission record.  No retained line contains a citation, so the wrapper carries no bibliography and no bibliography entry.  No retained line contains a figure environment, an image-inclusion command or drawing code, so no image is delivered and the package declares no asset dependency.  Editorial formatting only, in addition: an AMS \texttt{amsart} preamble that restates the article's own declarations and macros used by the retained lines, namely the environments \texttt{proposition} on separate counters, together with the standard packages those lines need.  The retained environments print in this document as Proposition 1 in order of appearance, whereas the article prints the same items with the numbering of its own full document.  The complete native source of the article as distributed in the arXiv e-print for this exact version is bundled unchanged as \texttt{sources/183671/source0.tex}.
\begin{equation}
\label{eq:convex-perspective}
    \bigl(\operatorname{cl}\widetilde h\bigr)(\mathbf v,y)
    \;=\;
    \begin{cases}
        \dfrac{\mathbf v^{\mathsf T}Q\mathbf v}{y}
        +\mathbf c^{\mathsf T}\mathbf v
        +dy, & y>0,\\[6pt]
        0, & y=0,
    \end{cases}
    \qquad \le 0,
\end{equation}

\begin{proposition}[Conic epigraph reformulation of HR]
\label{prop:conic-epi-equiv}
Let $Q\succeq 0$ and define
\begin{equation}
\label{eq:conic-set-S1}
S_{1}
=
\Bigl\{(\mathbf v,y)\in\mathbb R^{n}\times\mathbb R \ \Bigm|\ 
y\in[0,1],\ 
\bigl(\operatorname{cl}\widetilde h\bigr)(\mathbf v,y)\le 0,\ 
\mathbf x^{\ell}y\le \mathbf v\le \mathbf x^{u}y
\Bigr\},
\end{equation}
where $\bigl(\operatorname{cl}\widetilde h\bigr)$ is given by \eqref{eq:convex-perspective}. Define also
\begin{equation}
\label{eq:conic-set-S2}
\begin{aligned}
S_{2}^{\mathrm{conic}}
=
\Bigl\{(\mathbf v,y)\in\mathbb R^{n}\times\mathbb R \;\Bigm|\;&
y\in[0,1],\ \exists\, t\in\mathbb R,\ 
\mathbf v^{\mathsf T}Q\mathbf v \le t\,y,\\
& t + \mathbf c^{\mathsf T}\mathbf v + d\,y \le 0,\ 
t\ge 0,\ 
\mathbf x^{\ell}y\le \mathbf v\le \mathbf x^{u}y
\Bigr\}.
\end{aligned}
\end{equation}
Then $S_{1} = S_{2}^{\mathrm{conic}}$.
\end{proposition}

\begin{proof}
We prove this by mutual containment of the sets.

First, take $(\mathbf v,y)\in S_{1}$.
If $y=0$, the bounds $\mathbf x^{\ell}y\le \mathbf v\le \mathbf x^{u}y$ imply $\mathbf v=\mathbf 0$.
Then the choice $t=0$ satisfies $\mathbf v^{\mathsf T}Q\mathbf v = 0 \le t\,y=0$ and
$t+\mathbf c^{\mathsf T}\mathbf v + d\,y = 0$. Moreover, for any feasible $t$ we must have
$t+\mathbf c^{\mathsf T}\mathbf 0 + d\cdot 0 \le 0$ and $t\ge 0$, hence $t=0$ is forced.
Therefore $(\mathbf v,y)\in S_{2}^{\mathrm{conic}}$.
If $y>0$, define $t := \frac{\mathbf v^{\mathsf T}Q\mathbf v}{y}\ge 0$ (since $Q\succeq 0$).
Then $\mathbf v^{\mathsf T}Q\mathbf v \le t\,y$ holds at equality, and
$\bigl(\operatorname{cl}\widetilde h\bigr)(\mathbf v,y)\le 0$ implies
$t+\mathbf c^{\mathsf T}\mathbf v + d\,y \le 0$. Hence $(\mathbf v,y)\in S_{2}^{\mathrm{conic}}$.

Second, take $(\mathbf v,y)\in S_{2}^{\mathrm{conic}}$ with some $t\ge 0$ satisfying \eqref{eq:conic-set-S2}.
If $y=0$, the bounds imply $\mathbf v=\mathbf 0$, and then
$\bigl(\operatorname{cl}\widetilde h\bigr)(\mathbf 0,0)=0\le 0$, so $(\mathbf v,y)\in S_{1}$.
If $y>0$, the inequality $\mathbf v^{\mathsf T}Q\mathbf v \le t\,y$ implies
$\frac{\mathbf v^{\mathsf T}Q\mathbf v}{y}\le t$, and combining with
$t + \mathbf c^{\mathsf T}\mathbf v + d\,y \le 0$ yields
$\frac{\mathbf v^{\mathsf T}Q\mathbf v}{y} + \mathbf c^{\mathsf T}\mathbf v + d\,y \le 0$,
which is exactly $\bigl(\operatorname{cl}\widetilde h\bigr)(\mathbf v,y)\le 0$ for $y>0$.
Therefore $(\mathbf v,y)\in S_{1}$.

Thus $S_{1} = S_{2}^{\mathrm{conic}}$.
\end{proof}

\end{document}
